\section{Conclusiones}

El caso de la Universidad XUNA muestra que los problemas administrativos descritos al inicio 
(expedientes duplicados, información inconsistente, búsquedas lentas) no se resuelven simplemente centralizando los datos,
sino decidiendo con cuidado cómo se representan y cómo se accede a ellos. Esa fue la línea que guió todo el diseño propuesto en este informe.

Separar el registro en tres structs, en vez de un bloque único, salió de mirar qué tan seguido cambia cada dato y quién lo usa junto con qué, no de una preferencia de diseño. 

Lo mismo con el mapa de posiciones: no tiene sentido cargar 25\,000 expedientes completos en RAM si alcanza con guardar el carné, la posición en disco y el estado de cada uno.

El patrón ISAM terminó siendo lo que conecta arreglo, archivo secuencial y archivo indexado entre sí: 
la zona ordenada resuelve la búsqueda en $O(\log n)$, 
la lista de recién ingresados evita reordenar el archivo en cada matrícula nueva, 
y el offset calculado a partir del tamaño fijo del registro es lo que permite llegar directo a un registro en disco sin recorrer los demás.

En conjunto, estas decisiones muestran que el diseño de estructuras de datos no depende únicamente de que una solución funcione, 
sino de que su costo computacional se mantenga razonable a medida que la cantidad dato crece. 
Ese es el criterio que se aplicó en cada función documentada a lo largo del informe, y el que deberá mantenerse durante la implementación del sistema en el resto del semestre.
